\documentclass[12pt]{article}
% a4paper

%\usepackage[dvips]{epsfig}
%\usepackage{cours11, cours}
%\usepackage{fancyheadings}
%\usepackage{calc}
\usepackage{amsmath}
\usepackage{amssymb}
\usepackage{latexsym}
\usepackage{epsfig}
\usepackage{enumerate}
%\usepackage{supertabular}
%\usepackage{wrapfig}
%\usepackage{blackbrd}
%\usepackage{epic,eepic}
\usepackage{rotating}
\usepackage{multicol}
%\usepackage{multirow}
\usepackage{makeidx} % additional command see
\usepackage{rotating}
\usepackage{array}
\usepackage{tikz}
\usepackage{longtable}
%\usepackage{anyfontsize}
\usepackage{t1enc}
%\usepackage{amsmath,amsfonts}


%\usepackage[mtbold,mtplusscr]{mathtime}
% lucidacal,lucidascr,

%\usepackage{mathtimy}
%\usepackage{bm}
%\usepackage{avant}
%\usepackage{basker}
%\usepackage{bembo}
%\usepackage{bookman}
%\usepackage{chancery}
%\usepackage{garamond}
%\usepackage{helvet}
%\usepackage{newcent}
%\usepackage{palatino}
%\usepackage{times}
%\usepackage{pifont}



%\parindent=0pt


%\topmargin=0pt
%\headsep=18pt
%\footskip=45pt
%\mathsurround=1pt
%\evensidemargin=0pt
%\oddsidemargin=15pt

%\setlength{\textheight}{\baselineskip*41+\topskip}

\newcommand{\sectionline}{
   \nointerlineskip \vspace{\baselineskip}
   \hspace{\fill}\rule{0.9\linewidth}{1.7pt}\hspace{\fill}
   \par\nointerlineskip \vspace{\baselineskip}
   }
\newcommand\setTBstruts{\def\T{\rule{0pt}{2.6ex}}%
\def\B{\rule[-1.2ex]{0pt}{0pt}}}
\newcommand{\ans}[1]{\\{\bf ANSWER}: {#1}}
\newcommand{\Aut}{{\rm Aut}}
\newcommand{\Sym}{{\rm Sym}}
\newcommand{\sFix}{{\cal Fix}}
\newcommand{\sOrbits}{{\cal Orbits}}
\newcommand{\Stab}{{\rm Stab}}
\newcommand{\Fix}{{\rm Fix}}
\newcommand{\fix}{{\rm fix}}
\newcommand{\Orbits}{{\rm Orbits}}
\newcommand{\PG}{{\rm PG}}
\newcommand{\AG}{{\rm AG}}
\newcommand{\SQS}{{\rm SQS}}
\newcommand{\STS}{{\rm STS}}
\newcommand{\PSL}{{\rm PSL}}
\newcommand{\PGL}{{\rm PGL}}
\newcommand{\PSSL}{{\rm P\Sigma L}}
\newcommand{\PGGL}{{\rm P\Gamma L}}
\newcommand{\SL}{{\rm SL}}
\newcommand{\GL}{{\rm GL}}
\newcommand{\SSL}{{\rm \Sigma L}}
\newcommand{\GGL}{{\rm \Gamma L}}
\newcommand{\ASL}{{\rm ASL}}
\newcommand{\AGL}{{\rm AGL}}
\newcommand{\ASSL}{{\rm A\Sigma L}}
\newcommand{\AGGL}{{\rm A\Gamma L}}
\newcommand{\PSU}{{\rm PSU}}
\newcommand{\HS}{{\rm HS}}
\newcommand{\Hol}{{\rm Hol}}
\newcommand{\SO}{{\rm SO}}
\newcommand{\ASO}{{\rm ASO}}
\newcommand{\la}{\langle}
\newcommand{\ra}{\rangle}
\newcommand{\cA}{{\cal A}}
\newcommand{\cB}{{\cal B}}
\newcommand{\cC}{{\cal C}}
\newcommand{\cD}{{\cal D}}
\newcommand{\cE}{{\cal E}}
\newcommand{\cF}{{\cal F}}
\newcommand{\cG}{{\cal G}}
\newcommand{\cH}{{\cal H}}
\newcommand{\cI}{{\cal I}}
\newcommand{\cJ}{{\cal J}}
\newcommand{\cK}{{\cal K}}
\newcommand{\cL}{{\cal L}}
\newcommand{\cM}{{\cal M}}
\newcommand{\cN}{{\cal N}}
\newcommand{\cO}{{\cal O}}
\newcommand{\cP}{{\cal P}}
\newcommand{\cQ}{{\cal Q}}
\newcommand{\cR}{{\cal R}}
\newcommand{\cS}{{\cal S}}
\newcommand{\cT}{{\cal T}}
\newcommand{\cU}{{\cal U}}
\newcommand{\cV}{{\cal V}}
\newcommand{\cW}{{\cal W}}
\newcommand{\cX}{{\cal X}}
\newcommand{\cY}{{\cal Y}}
\newcommand{\cZ}{{\cal Z}}
\newcommand{\rmA}{{\rm A}}
\newcommand{\rmB}{{\rm B}}
\newcommand{\rmC}{{\rm C}}
\newcommand{\rmD}{{\rm D}}
\newcommand{\rmE}{{\rm E}}
\newcommand{\rmF}{{\rm F}}
\newcommand{\rmG}{{\rm G}}
\newcommand{\rmH}{{\rm H}}
\newcommand{\rmI}{{\rm I}}
\newcommand{\rmJ}{{\rm J}}
\newcommand{\rmK}{{\rm K}}
\newcommand{\rmL}{{\rm L}}
\newcommand{\rmM}{{\rm M}}
\newcommand{\rmN}{{\rm N}}
\newcommand{\rmO}{{\rm O}}
\newcommand{\rmP}{{\rm P}}
\newcommand{\rmQ}{{\rm Q}}
\newcommand{\rmR}{{\rm R}}
\newcommand{\rmS}{{\rm S}}
\newcommand{\rmT}{{\rm T}}
\newcommand{\rmU}{{\rm U}}
\newcommand{\rmV}{{\rm V}}
\newcommand{\rmW}{{\rm W}}
\newcommand{\rmX}{{\rm X}}
\newcommand{\rmY}{{\rm Y}}
\newcommand{\rmZ}{{\rm Z}}
\newcommand{\bA}{{\bf A}}
\newcommand{\bB}{{\bf B}}
\newcommand{\bC}{{\bf C}}
\newcommand{\bD}{{\bf D}}
\newcommand{\bE}{{\bf E}}
\newcommand{\bF}{{\bf F}}
\newcommand{\bG}{{\bf G}}
\newcommand{\bH}{{\bf H}}
\newcommand{\bI}{{\bf I}}
\newcommand{\bJ}{{\bf J}}
\newcommand{\bK}{{\bf K}}
\newcommand{\bL}{{\bf L}}
\newcommand{\bM}{{\bf M}}
\newcommand{\bN}{{\bf N}}
\newcommand{\bO}{{\bf O}}
\newcommand{\bP}{{\bf P}}
\newcommand{\bQ}{{\bf Q}}
\newcommand{\bR}{{\bf R}}
\newcommand{\bS}{{\bf S}}
\newcommand{\bT}{{\bf T}}
\newcommand{\bU}{{\bf U}}
\newcommand{\bV}{{\bf V}}
\newcommand{\bW}{{\bf W}}
\newcommand{\bX}{{\bf X}}
\newcommand{\bY}{{\bf Y}}
\newcommand{\bZ}{{\bf Z}}
\newcommand{\sA}{{\cal A}}
\newcommand{\sB}{{\cal B}}
\newcommand{\sC}{{\cal C}}
\newcommand{\sD}{{\cal D}}
\newcommand{\sE}{{\cal E}}
\newcommand{\sF}{{\cal F}}
\newcommand{\sG}{{\cal G}}
\newcommand{\sH}{{\cal H}}
\newcommand{\sI}{{\cal I}}
\newcommand{\sJ}{{\cal J}}
\newcommand{\sK}{{\cal K}}
\newcommand{\sL}{{\cal L}}
\newcommand{\sM}{{\cal M}}
\newcommand{\sN}{{\cal N}}
\newcommand{\sO}{{\cal O}}
\newcommand{\sP}{{\cal P}}
\newcommand{\sQ}{{\cal Q}}
\newcommand{\sR}{{\cal R}}
\newcommand{\sS}{{\cal S}}
\newcommand{\sT}{{\cal T}}
\newcommand{\sU}{{\cal U}}
\newcommand{\sV}{{\cal V}}
\newcommand{\sW}{{\cal W}}
\newcommand{\sX}{{\cal X}}
\newcommand{\sY}{{\cal Y}}
\newcommand{\sZ}{{\cal Z}}
\newcommand{\frakA}{{\mathfrak A}}
\newcommand{\frakB}{{\mathfrak B}}
\newcommand{\frakC}{{\mathfrak C}}
\newcommand{\frakD}{{\mathfrak D}}
\newcommand{\frakE}{{\mathfrak E}}
\newcommand{\frakF}{{\mathfrak F}}
\newcommand{\frakG}{{\mathfrak G}}
\newcommand{\frakH}{{\mathfrak H}}
\newcommand{\frakI}{{\mathfrak I}}
\newcommand{\frakJ}{{\mathfrak J}}
\newcommand{\frakK}{{\mathfrak K}}
\newcommand{\frakL}{{\mathfrak L}}
\newcommand{\frakM}{{\mathfrak M}}
\newcommand{\frakN}{{\mathfrak N}}
\newcommand{\frakO}{{\mathfrak O}}
\newcommand{\frakP}{{\mathfrak P}}
\newcommand{\frakQ}{{\mathfrak Q}}
\newcommand{\frakR}{{\mathfrak R}}
\newcommand{\frakS}{{\mathfrak S}}
\newcommand{\frakT}{{\mathfrak T}}
\newcommand{\frakU}{{\mathfrak U}}
\newcommand{\frakV}{{\mathfrak V}}
\newcommand{\frakW}{{\mathfrak W}}
\newcommand{\frakX}{{\mathfrak X}}
\newcommand{\frakY}{{\mathfrak Y}}
\newcommand{\frakZ}{{\mathfrak Z}}
\newcommand{\fraka}{{\mathfrak a}}
\newcommand{\frakb}{{\mathfrak b}}
\newcommand{\frakc}{{\mathfrak c}}
\newcommand{\frakd}{{\mathfrak d}}
\newcommand{\frake}{{\mathfrak e}}
\newcommand{\frakf}{{\mathfrak f}}
\newcommand{\frakg}{{\mathfrak g}}
\newcommand{\frakh}{{\mathfrak h}}
\newcommand{\fraki}{{\mathfrak i}}
\newcommand{\frakj}{{\mathfrak j}}
\newcommand{\frakk}{{\mathfrak k}}
\newcommand{\frakl}{{\mathfrak l}}
\newcommand{\frakm}{{\mathfrak m}}
\newcommand{\frakn}{{\mathfrak n}}
\newcommand{\frako}{{\mathfrak o}}
\newcommand{\frakp}{{\mathfrak p}}
\newcommand{\frakq}{{\mathfrak q}}
\newcommand{\frakr}{{\mathfrak r}}
\newcommand{\fraks}{{\mathfrak s}}
\newcommand{\frakt}{{\mathfrak t}}
\newcommand{\fraku}{{\mathfrak u}}
\newcommand{\frakv}{{\mathfrak v}}
\newcommand{\frakw}{{\mathfrak w}}
\newcommand{\frakx}{{\mathfrak x}}
\newcommand{\fraky}{{\mathfrak y}}
\newcommand{\frakz}{{\mathfrak z}}
\newcommand{\Tetra}{{\mathfrak Tetra}}
\newcommand{\Cube}{{\mathfrak Cube}}
\newcommand{\Octa}{{\mathfrak Octa}}
\newcommand{\Dode}{{\mathfrak Dode}}
\newcommand{\Ico}{{\mathfrak Ico}}
\newcommand{\bbF}{{\mathbb F}}
\newcommand{\bbQ}{{\mathbb Q}}
\newcommand{\bbC}{{\mathbb C}}
\newcommand{\bbR}{{\mathbb R}}



%\makeindex

\begin{document} 
\setTBstruts

\bibliographystyle{plain}
%\large

{\allowdisplaybreaks%




%\makeindex

%\renewcommand{\labelenumi}{(\roman{enumi})}

\title{Varieties of degree 3 in PG(2,4)}
\author{Orbiter}%end author
%\date{}
\maketitle%
%\pagenumbering{roman}
%\thispagestyle{empty}
%\input et.tex%
%\thispagestyle{empty}%\phantom{page2}%\clearpage%
%\addcontentsline{toc}{chapter}{Inhaltsverzeichnis}%
%\tableofcontents
%\listofsymbols
%\pagenumbering{arabic}
%\pagenumbering{roman}



\small
\arraycolsep=2pt
\parindent=0pt
$q = 4$\\
$n = 2$\\
degree of poly $ = 3$\\
\clearpage

\section*{The Varieties of degree $3$ in $PG(2, 4)$, summary}
\subsection*{Objects with 13 Points}
There are 1 objects with 13 Points: \\
orbit : rep : go : poly : Pts\\
4 : ( 11 ) : 288 : 11=$X^3 + Y^3$ : "2, 3, 4, 5, 6, 11, 12, 14, 15, 16, 18, 19, 20"\\
\subsection*{Objects with 12 Points}
There are 1 objects with 12 Points: \\
orbit : rep : go : poly : Pts\\
2 : ( 9 ) : 54 : 9=$XYZ$ : "0, 1, 2, 4, 5, 6, 7, 8, 9, 10, 13, 17"\\
\subsection*{Objects with 10 Points}
There are 1 objects with 10 Points: \\
orbit : rep : go : poly : Pts\\
31 : ( 87743 ) : 10 : 87743=$X^3 + 2Y^3 + 2Z^3 + X^2Y + X^2Z + XYZ$ : "3, 5, 8, 10, 13, 14, 15, 17, 18, 20"\\
\subsection*{Objects with 9 Points}
There are 4 objects with 9 Points: \\
orbit : rep : go : poly : Pts\\
1 : ( 3 ) : 144 : 3=$X^2Y$ : "0, 1, 2, 7, 8, 9, 10, 13, 17"\\
6 : ( 18 ) : 216 : 18=$X^3 + Y^3 + Z^3$ : "4, 5, 6, 7, 8, 9, 10, 13, 17"\\
9 : ( 24 ) : 27 : 24=$3X^3 + 2Y^3 + Z^3$ : "3, 11, 12, 14, 15, 16, 18, 19, 20"\\
15 : ( 175 ) : 12 : 175=$Y^3 + Z^3 + X^2Y + X^2Z$ : "0, 3, 4, 7, 10, 11, 12, 13, 17"\\
\subsection*{Objects with 8 Points}
There are 2 objects with 8 Points: \\
orbit : rep : go : poly : Pts\\
19 : ( 87383 ) : 6 : 87383=$X^3 + XYZ$ : "1, 2, 3, 10, 13, 16, 17, 19"\\
24 : ( 87466 ) : 2 : 87466=$Y^3 + Z^3 + X^2Y + XYZ$ : "0, 3, 4, 10, 13, 14, 17, 18"\\
\subsection*{Objects with 7 Points}
There are 2 objects with 7 Points: \\
orbit : rep : go : poly : Pts\\
12 : ( 50 ) : 6 : 50=$2X^3 + Y^3 + Z^3 + X^2Y$ : "5, 10, 12, 13, 14, 17, 19"\\
13 : ( 51 ) : 6 : 51=$3X^3 + Y^3 + Z^3 + X^2Y$ : "6, 10, 11, 13, 16, 17, 18"\\
\subsection*{Objects with 6 Points}
There are 6 objects with 6 Points: \\
orbit : rep : go : poly : Pts\\
21 : ( 87404 ) : 6 : 87404=$2X^3 + Y^3 + Z^3 + XYZ$ : "10, 11, 13, 14, 17, 20"\\
22 : ( 87405 ) : 6 : 87405=$3X^3 + Y^3 + Z^3 + XYZ$ : "10, 12, 13, 15, 17, 18"\\
25 : ( 87468 ) : 3 : 87468=$2X^3 + Y^3 + Z^3 + X^2Y + XYZ$ : "5, 10, 13, 15, 16, 17"\\
26 : ( 87469 ) : 3 : 87469=$3X^3 + Y^3 + Z^3 + X^2Y + XYZ$ : "6, 10, 13, 17, 19, 20"\\
28 : ( 87471 ) : 2 : 87471=$X^3 + 2Y^3 + Z^3 + X^2Y + XYZ$ : "5, 7, 8, 9, 15, 16"\\
32 : ( 91990 ) : 30 : 91990=$2X^2Z + Y^2Z + XYZ$ : "0, 1, 2, 4, 5, 6"\\
\subsection*{Objects with 5 Points}
There are 4 objects with 5 Points: \\
orbit : rep : go : poly : Pts\\
0 : ( 0 ) : 2880 : 0=$X^3$ : "1, 2, 10, 13, 17"\\
10 : ( 34 ) : 96 : 34=$2X^3 + Y^3 + X^2Y$ : "2, 5, 11, 16, 18"\\
11 : ( 44 ) : 3 : 44=$Z^3 + X^2Y$ : "0, 1, 3, 15, 20"\\
16 : ( 176 ) : 4 : 176=$X^3 + Y^3 + Z^3 + X^2Y + X^2Z$ : "10, 13, 16, 17, 19"\\
\subsection*{Objects with 4 Points}
There are 4 objects with 4 Points: \\
orbit : rep : go : poly : Pts\\
3 : ( 10 ) : 6 : 10=$X^3 + Y^3 + Z^3 + X^2Y + X^2Z + XY^2 + Y^2Z + XZ^2 + YZ^2 + XYZ$ : "3, 4, 7, 10"\\
20 : ( 87388 ) : 3 : 87388=$2X^3 + Y^3 + XYZ$ : "2, 12, 15, 18"\\
27 : ( 87470 ) : 2 : 87470=$2Y^3 + Z^3 + X^2Y + XYZ$ : "0, 6, 19, 20"\\
30 : ( 87474 ) : 2 : 87474=$3Y^3 + Z^3 + X^2Y + XYZ$ : "0, 5, 15, 16"\\
\subsection*{Objects with 3 Points}
There are 4 objects with 3 Points: \\
orbit : rep : go : poly : Pts\\
7 : ( 19 ) : 18 : 19=$2X^3 + Y^3 + Z^3$ : "10, 13, 17"\\
8 : ( 20 ) : 18 : 20=$3X^3 + Y^3 + Z^3$ : "10, 13, 17"\\
17 : ( 180 ) : 9 : 180=$X^3 + 2Y^3 + Z^3 + X^2Y + X^2Z$ : "5, 15, 18"\\
18 : ( 181 ) : 9 : 181=$2X^3 + 2Y^3 + Z^3 + X^2Y + X^2Z$ : "8, 16, 19"\\
\subsection*{Objects with 2 Points}
There are 1 objects with 2 Points: \\
orbit : rep : go : poly : Pts\\
29 : ( 87472 ) : 2 : 87472=$2X^3 + 2Y^3 + Z^3 + X^2Y + XYZ$ : "11, 12"\\
\subsection*{Objects with 1 Points}
There are 2 objects with 1 Points: \\
orbit : rep : go : poly : Pts\\
5 : ( 12 ) : 144 : 12=$2X^3 + Y^3$ : "2"\\
14 : ( 55 ) : 24 : 55=$3X^3 + 2Y^3 + Z^3 + X^2Y$ : "4"\\
\subsection*{Objects with 0 Points}
There are 1 objects with 0 Points: \\
orbit : rep : go : poly : Pts\\
23 : ( 87409 ) : 63 : 87409=$3X^3 + 2Y^3 + Z^3 + XYZ$ : ""\\


%\bibliographystyle{gerplain}% wird oben eingestellt
%\addcontentsline{toc}{section}{References}
%\bibliography{../MY_BIBLIOGRAPHY/anton}
% ACHTUNG: nicht vergessen:
% die Zeile
%\addcontentsline{toc}{chapter}{Literaturverzeichnis}
% muss per Hand in d.bbl eingefuegt werden !
% nach \begin{thebibliography}{100}

%\begin{theindex}

%\clearpage
%\addcontentsline{toc}{chapter}{Index}
%\input{apd.ind}

%\printindex
%\end{theindex}

}% allowdisplaybreaks

\end{document}


